\documentclass[12pt]{article}
\usepackage[spanish]{babel}
\usepackage[utf8]{inputenc}
\usepackage[T1]{fontenc}
\usepackage{graphicx}
\usepackage{geometry}
\usepackage{float}
\usepackage{subcaption}
\usepackage[most]{tcolorbox}
\usepackage{hyperref}
\usepackage{longtable}
\usepackage{listings}
\usepackage{xcolor}
\geometry{a4paper, margin=2.5cm}

\title{Manual de Usuario}
\author{Cynthia Castaño Juan}
\date{12 de abril de 2026}

\begin{document}

\maketitle
\tableofcontents
\newpage

\section{Anteproyecto de Trabajo Fin de Grado: Aplicación móvil de eventos emergentes en Sevilla, ``Sevillaneando''}

\subsection{Descripción}
El propósito de la aplicación es promover eventos culturales y de ocio poco conocidos o emergentes, como talleres, exposiciones, monumentos y conciertos. No obstante, también ofrecerá información sobre eventos convencionales.

Estos eventos estarán listados por categorías, con información detallada de cada uno de ellos: mapa, reseñas, imágenes, descripción, precio, entre otros. Además, se permitirá a los usuarios añadir nuevos eventos (sujetos a revisión y aprobación por parte de un moderador) y compartir fotos en tiempo real transient los eventos estén activos, fomentando la visibilidad y la participación de la comunidad.

Por último, la aplicación ofrecerá recomendaciones personalizadas y notificaciones push para mantener informados a los usuarios en función de su ubicación e intereses.

\subsection{Objetivos}
\begin{itemize}
    \item Listar información detallada sobre eventos culturales y de ocio en Sevilla.
    \item Permitir a los usuarios añadir eventos para promover iniciativas culturales o negocios emergentes, los cuales serán revisados y aprobados por un moderador antes de su publicación.
    \item Integrar un mapa para localizar los eventos.
    \item Habilitar la subida de imágenes en vivo durante los eventos.
    \item Ofrecer reseñas y valoraciones de los eventos.
    \item Incorporar recomendaciones personalizadas basadas en preferencias y localización.
    \item Enviar notificaciones push para avisar de eventos cercanos o recomendados.
\end{itemize}

\subsection{Tecnologías iniciales}

\subsubsection{Frontend móvil (aplicación)}
\textbf{Alternativas:}
\begin{itemize}
    \item \textbf{React Native (JavaScript)}: framework desarrollado por Facebook, ampliamente utilizado, que permite crear aplicaciones nativas para iOS y Android con una única base de código.
    \item \textbf{Flutter (Dart)}: framework de Google de alto rendimiento, con widgets propios y desarrollo multiplataforma.
    \item \textbf{Nativo (Kotlin para Android / Swift para iOS)}: ofrece el máximo rendimiento, pero requiere desarrollo duplicado para cada plataforma.
\end{itemize}

\textbf{Decisión final:} se opta por \textbf{React Native}, ya que permite una única base de código para iOS y Android, dispone de una amplia comunidad, numerosas librerías (mapas, geolocalización, notificaciones), abundante documentación y una mayor adopción en comparación con Flutter en el contexto del proyecto.

\subsubsection{Backend / API}
\textbf{Alternativas:}
\begin{itemize}
    \item \textbf{Node.js con Express/NestJS}: rápido, escalable y con integración directa con el frontend en React Native.
    \item \textbf{Python con Django REST Framework o FastAPI}: muy adecuado para APIs REST y con fuerte soporte en inteligencia artificial.
    \item \textbf{Serverless (Firebase Functions, AWS Lambda)}: no requiere servidor propio y ofrece escalado automático.
\end{itemize}

\textbf{Decisión final:} se opta por \textbf{Node.js con NestJS}, ya que mantiene coherencia con el lenguaje del frontend (JavaScript/TypeScript), permite desarrollar APIs REST de forma estructurada y facilita la implementación de autenticación y comunicación en tiempo real gracias a su arquitectura basada en eventos.

\subsubsection{Base de datos}
\textbf{Alternativas:}
\begin{itemize}
    \item \textbf{PostgreSQL}: base de datos relacional muy estable con soporte geoespacial mediante PostGIS.
    \item \textbf{MySQL/MariaDB}: soluciones relacionales ampliamente utilizadas, pero con menor soporte geoespacial avanzado.
    \item \textbf{MongoDB}: base de datos NoSQL flexible para datos semiestructurados.
\end{itemize}

\textbf{Decisión final:} se opta por \textbf{PostgreSQL con PostGIS}, ya que los eventos incluyen ubicaciones geográficas y requieren consultas por proximidad. PostGIS proporciona funciones geoespaciales nativas como distancias, áreas y búsqueda de puntos cercanos.

\subsubsection{Autenticación}
\textbf{Alternativas:}
\begin{itemize}
    \item \textbf{Firebase Auth}: integración rápida, soporte para email y proveedores externos como Google.
    \item \textbf{Auth0}: solución potente y configurable orientada a entornos empresariales.
    \item \textbf{JWT propio}: implementación personalizada directamente en el backend.
\end{itemize}

\textbf{Decisión final:} se opta por \textbf{Firebase Auth}, ya que reduce significativamente el tiempo de desarrollo, gestiona autenticación de forma segura y soporta proveedores sociales sin necesidad de infraestructura propia.

\subsubsection{Almacenamiento de imágenes}
\textbf{Alternativas:}
\begin{itemize}
    \item Almacenamiento local en el backend.
    \item \textbf{Firebase Storage}: integración directa con Firebase Auth y SDKs para React Native.
    \item \textbf{AWS S3}: estándar industrial altamente escalable y flexible.
    \item \textbf{Supabase Storage}: alternativa open source con modelo similar a Firebase.
\end{itemize}

\textbf{Decisión final:} se opta por \textbf{Firebase Storage}, debido a su integración directa con Firebase Auth, facilidad de uso y soporte nativo en React Native, lo que permite controlar fácilmente permisos de subida y acceso a imágenes.

\subsubsection{Mapas}
\textbf{Alternativas evaluadas:}
\begin{itemize}
    \item \textbf{Google Maps SDK}: solución muy completa y ampliamente utilizada, aunque con costes asociados por uso.
    \item \textbf{Mapbox}: plataforma altamente personalizable con un plan gratuito limitado, aunque con licencia no completamente abierta.
    \item \textbf{OpenStreetMap con react-native-maps}: solución libre y de código abierto que permite mapas interactivos sin costes de licencia.
\end{itemize}

\textbf{Decisión final:} se opta por \textbf{OpenStreetMap integrado mediante react-native-maps}, ya que permite representar mapas interactivos sin costes asociados.

Cada evento mostrará un mapa embebido con su ubicación exacta mediante un marcador basado en coordenadas GPS reales de la ciudad de Sevilla.

Además, desde la ficha del evento se permitirá iniciar una ruta hacia la ubicación seleccionada mediante aplicaciones externas como Google Maps o Apple Maps, proporcionando una experiencia de usuario completa sin necesidad de implementar un sistema de navegación propio.

\subsubsection{Moderación automática de contenido (imágenes en vivo)}
\textbf{Alternativas:}
\begin{itemize}
    \item \textbf{Google Vision SafeSearch API}: clasificación de contenido sensible (desnudez, violencia), aunque con coste asociado.
    \item \textbf{AWS Rekognition}: solución similar con capacidades avanzadas de análisis de imágenes.
    \item \textbf{Modelos propios (NSFW.js / TensorFlow.js)}: solución gratuita con mayor control, pero requiere ajuste y mantenimiento.
\end{itemize}

\textbf{Decisión final:} se opta por un \textbf{modelo propio basado en NSFW.js con TensorFlow.js}, ya que no tiene coste, permite control total sobre los datos y se integra fácilmente con el backend en Node.js para analizar imágenes antes de su publicación o almacenamiento.


\end{document}
